\documentclass[11pt]{article}
\usepackage{docstyle}
% Lesson 2 lost-mark points (from the material plus teaching observation), numbered once
\setcounter{lostmark}{0}
\lostmark{noy}{Stopping at $x$: a \tturn is a point, so find $y$ from $f(x)$ as well.}
\lostmark{nojust}{Writing ``maximum'' with no reason. Justify with $f''(x)$ or the gradient on each side.}
\lostmark{interval}{Increasing or decreasing: giving the \tturns instead of an interval, or the inequality the wrong way round.}
\lostmark{twovars}{Differentiating a formula that still has two variables. Use the constraint to get one variable first.}
\lostmark{wrongq}{Answering the wrong quantity: the question asks for the maximum area, not the $x$ that gives it.}
\lostmark{context}{No units or context in the final sentence (m, s, \$, people per hour).}

\newlength\figunit
\newcommand\ddxx[1]{\dfrac{\mathrm{d}^{2}#1}{\mathrm{d}x^{2}}}
\begin{document}
\handouthead{Lesson 2 \textperiodcentered\ Turning Points and Optimisation}{NCEA Level 2 Mathematics \textperiodcentered\ AS 91262 Apply calculus methods in solving problems \textperiodcentered\ Handout}

\section{The one idea}\label{idea}
\begin{minipage}[c]{0.54\linewidth}
At the top of a hill or the bottom of a valley the \ttangent is flat.
\begin{keybox}
At a \tturn the gradient is zero: $\mathbf{f'(x)=0}$.\\
Maximum or minimum? Look at the gradient on each side.
\end{keybox}
On the left of a \kw{maximum} the curve goes up ($+$), on the right it goes down ($-$).
A \kw{minimum} is the other way round: $-$ then $+$.
\end{minipage}\hfill
\begin{minipage}[c]{0.42\linewidth}\centering
\setlength\figunit{11mm}
% f(x) = x^3 - 6x^2 + 9x + 1: maximum (1,5), minimum (3,1).  Flat tangents drawn at both;
% gradient signs from f'(x) = 3(x-1)(x-3): f'(0)=9>0, f'(2)=-3<0, f'(4)=9>0.  Caller sets \figunit.
\begin{tikzpicture}[x=\figunit, y=0.55\figunit, line cap=round]
  \draw[-{Stealth[length=2mm]}, line width=0.6pt] (-0.6,0) -- (5.6,0) node[right, inner sep=1pt] {$x$};
  \draw[-{Stealth[length=2mm]}, line width=0.6pt] (0,-0.5) -- (0,7.3) node[above, inner sep=1pt] {$y$};
  \begin{scope}
    \clip (-0.6,-0.5) rectangle (4.8,7.2);
    \curve{-0.25}{4.45}{\x*\x*\x-6*\x*\x+9*\x+1}
  \end{scope}
  \draw[line width=0.8pt, dashed] (0.4,5) -- (1.6,5);
  \draw[line width=0.8pt, dashed] (2.4,1) -- (3.6,1);
  \fill (1,5) circle (2.2pt); \fill (3,1) circle (2.2pt);
  \node[above, inner sep=2pt] at (1,5.1) {maximum};
  % 'minimum' right of the flat tangent: at x = 3.7 the curve is at 2.8, above the label
  \node[anchor=west, inner sep=2pt] at (3.65,1) {minimum};
  % gradient signs, placed beside the curve
  \node at (0.6,2.2) {\textbf{$+$}};   % below the rising curve (f(0.6) = 4.5), clear of the y-axis
  \node at (2.35,4.3) {\textbf{$-$}};
  \node at (4.35,3.1) {\textbf{$+$}};
\end{tikzpicture}

\end{minipage}

\section{Finding turning points}\label{find}
\subsection{Write it in this order}
\begin{enumerate}[leftmargin=7mm, itemsep=1pt]
\item Differentiate: $f'(x)$.
\item Set $f'(x)=0$ and solve for $x$ (factorise, formula, or \textsf{EQUA}).
\item Find each $y$ from $f(x)$, not from $f'(x)$.
\item Decide maximum or minimum, and say why (section 3).
\end{enumerate}

\begin{example}{Worked example 1}
Find the \tturns of $f(x)=x^{3}-6x^{2}+9x+1$ and state their nature.
\begin{steps}
\item $f'(x)=3x^{2}-12x+9=3(x^{2}-4x+3)=3(x-1)(x-3)$
\item $f'(x)=0$ when $x=1$ or $x=3$.
\item $f(1)=1-6+9+1=5$\quad and\quad $f(3)=27-54+27+1=1$.
\item $f''(x)=6x-12$:\quad $f''(1)=-6<0$, so $\mathbf{(1,\,5)}$ \textbf{is a maximum};\quad $f''(3)=6>0$, so $\mathbf{(3,\,1)}$ \textbf{is a minimum}.
\end{steps}
\end{example}

\section{Maximum or minimum? Two tests}\label{nature}
\begin{tabular}{@{}p{0.23\linewidth}p{0.36\linewidth}p{0.33\linewidth}@{}}
\toprule
test & what you see & conclusion\\
\midrule
\kw{gradient each side} & $f'$ goes $+\ \to\ 0\ \to\ -$ & maximum \\
 & $f'$ goes $-\ \to\ 0\ \to\ +$ & minimum \\
\kw{second derivative} & $f''(x)<0$ at the point & maximum (curve bends down) \\
 & $f''(x)>0$ at the point & minimum (curve bends up) \\
\bottomrule
\end{tabular}

\medskip
$f''(x)$ means ``differentiate twice''; it is also written $\ddxx{y}$.
For the gradient test, choose $x$-values just either side and \emph{show} the substitution:
e.g.\ $f'(0)=9>0$ and $f'(2)=-3<0$ around $x=1$.

\textbf{Sentence frame:} ``$f''(1)=-6<0$, so the \tturn at $(1,\,5)$ is a maximum.''

\section{Increasing and decreasing}\label{incdec}
\begin{keybox}
\kw{Increasing} where $f'(x)>0$.\qquad \kw{Decreasing} where $f'(x)<0$.
\end{keybox}
Find where $f'(x)=0$ first; these $x$-values split the line into intervals. Then test each interval.
If $f'(x)$ is an up-parabola, it is \kw{negative between its roots}.

\begin{example}{Worked example 2}
For $f(x)=x^{3}-6x^{2}+9x+1$, find where $f$ is decreasing.
\begin{steps}
\item $f'(x)=3(x-1)(x-3)$, zero at $x=1$ and $x=3$.
\item Up-parabola, so negative between the roots: $\mathbf{1<x<3}$.
\end{steps}
Increasing for $x<1$ or $x>3$. Matches the picture: down between the maximum and the minimum.
\end{example}

\begin{example}{Worked example 3}
For which $x$ is $y=5+4x-x^{2}$ increasing?
\begin{steps}
\item $\dydx=4-2x$
\item $4-2x>0$ gives $2x<4$, so $\mathbf{x<2}$.
\end{steps}
\end{example}

\section{Optimisation: the biggest or smallest value}\label{opt}
\subsection{Write it in this order}
\begin{enumerate}[leftmargin=7mm, itemsep=1pt]
\item Write the quantity as a formula in \kw{one variable} (use the constraint to remove the other).
\item Differentiate.
\item Set the \tderiv $=0$ and solve.
\item Justify maximum or minimum ($f''$ or the gradient each side).
\item Answer the question asked, with units.
\end{enumerate}

\begin{example}{Worked example 4 \textperiodcentered\ formula given}
A ball's height is $h=2+19.6t-4.9t^{2}$ metres after $t$ seconds. Find the maximum height.
\begin{steps}
\item $\dd{h}{t}=19.6-9.8t=0$, so $t=2$ s.
\item $h=2+19.6(2)-4.9(2)^{2}=2+39.2-19.6=21.6$
\item $\dfrac{\mathrm{d}^{2}h}{\mathrm{d}t^{2}}=-9.8<0$, so this is a maximum. \textbf{Maximum height $21.6$ m.}
\end{steps}
\end{example}

\begin{example}{Worked example 5 \textperiodcentered\ build the formula (Merit/Excellence)}
\begin{minipage}[t]{0.60\linewidth}
A farmer has $60$ m of fence for a rectangular pen against a wall (three sides fenced). Find the maximum area.
\begin{steps}
\item Sides $x$, $x$ and $60-2x$:\quad $A=x(60-2x)=60x-2x^{2}$
\item $\dd{A}{x}=60-4x=0$, so $x=15$.
\item $\dfrac{\mathrm{d}^{2}A}{\mathrm{d}x^{2}}=-4<0$, so maximum.
\item $A=15\times30=\mathbf{450\ m^{2}}$ (pen $15$ m by $30$ m).
\end{steps}
\end{minipage}\hfill
\begin{minipage}[t]{0.36\linewidth}\centering\vspace{2mm}
\setlength\figunit{8mm}
% Rectangular pen against a wall: two sides x, one side 60 - 2x.  Caller sets \figunit.
\begin{tikzpicture}[x=\figunit, y=\figunit, line cap=round]
  \fill[pattern=north east lines, pattern color=gray] (-0.4,2.2) rectangle (5.4,2.5);
  \draw[line width=0.8pt] (-0.4,2.2) -- (5.4,2.2);
  \node[above, fill=white, inner sep=1pt] at (2.5,2.5) {wall};
  \draw[line width=1.2pt] (0,2.2) -- (0,0) -- (5,0) -- (5,2.2);
  \node[left] at (0,1.1) {$x$};
  \node[right] at (5,1.1) {$x$};
  \node[below] at (2.5,0) {$60-2x$};
\end{tikzpicture}

\end{minipage}
\end{example}

\section{What each grade looks like here}\label{grades}
\begin{tabular}{@{}p{0.14\linewidth}p{0.80\linewidth}@{}}
\toprule
Achieved & find a \tturn, or the maximum of a given formula, with the method shown.\\
Merit & link steps: justify the nature, find an interval, find a constant from a \tturn, interpret in context.\\
Excellence & build the model from words and a diagram, optimise it, justify, and answer in context; or show a general result with letters.\\
\bottomrule
\end{tabular}

\section{Checking with the fx-9750GIII}\label{calc}
\begin{itemize}[leftmargin=5mm, itemsep=1pt]
\item \textsf{GRAPH}: enter $Y1$, \textsf{DRAW}, then \textsf{G-SOLVE} (SHIFT F5) $\to$ \textsf{MAX} or \textsf{MIN}.
Use it to check a \tturn, never as the method.
\item \textsf{EQUA} $\to$ \textsf{POLY} solves $f'(x)=0$ when $f'$ is a quadratic.
\end{itemize}

\section{Ways to lose marks}\label{lose}
From the material plus teaching observation.
\begin{enumerate}[leftmargin=7mm, itemsep=2pt]
\item[\textbf{\lmnum{noy}}] \lmtxt{noy}
\item[\textbf{\lmnum{nojust}}] \lmtxt{nojust}
\item[\textbf{\lmnum{interval}}] \lmtxt{interval}
\item[\textbf{\lmnum{twovars}}] \lmtxt{twovars}
\item[\textbf{\lmnum{wrongq}}] \lmtxt{wrongq}
\item[\textbf{\lmnum{context}}] \lmtxt{context}
\end{enumerate}
\end{document}
